\section{One table, two searching questions}
The first reading of the whole formulary asks what makes a guest fit for communion. Gregory's inference from the expulsion locates the feast in the present Church; his garment is charity. Augustine supplies the accompanying logic of pardon, divine help and purification. The chief danger is to mistake entry for a completed life of love. Its characteristic response is to let received mercy become concrete charity, including toward an enemy or an inconvenient person in need.

The second asks what provision means when the faithful still hunger. Chrysostom refuses to make Paul's contentment cancel the Philippians' gift; Augustine refuses to let the prosperous wrongdoer define the psalm's good. Thomas places present supply within the distinction between the Church's banquet and heaven. The chief danger is to confuse blessing with possessions or, in reaction, to dismiss bodily need. Its characteristic response is to share distress and necessities while hoping for a fullness no gift can purchase.

The two readings strengthen each other at their most vulnerable points. The garment prevents abundance from becoming a comfortable entitlement. Paul's commendation of help prevents charity from becoming merely an inward condition. The psalm's forgiveness prevents moral seriousness from becoming self-salvation, while Isaiah's destruction of death prevents present suffering from becoming the last word. Neither reading requires every appointment to say the same thing. Their force comes from the different work done by pardon, call, gift, judgment, nourishment and hope.

With the short Gospel, the garment-centered argument needs its explicit qualification: the assembled guests are proclaimed, their inspection is not. With either Communion option, the final good remains more than immediate advantage. The psalm speaks of seeking the Lord; John speaks of seeing him. In both, the Mass's movement turns reception into a life capable of the communion it receives.
